\documentclass[12pt,a4paper]{article}

\usepackage[utf8]{inputenc}
\usepackage[T1]{fontenc}
\usepackage[romanian]{babel}
\usepackage{geometry}
\usepackage{booktabs}
\usepackage{array}
\usepackage{longtable}
\usepackage{hyperref}
\usepackage{xcolor}
\usepackage{listings}

\geometry{margin=2.5cm}

\definecolor{codebackground}{rgb}{0.95,0.95,0.95}
\lstdefinestyle{codestyle}{
    backgroundcolor=\color{codebackground},
    basicstyle=\ttfamily\small,
    breaklines=true,
    frame=single,
    keepspaces=true,
    showstringspaces=false,
    tabsize=2
}
\lstset{style=codestyle}

\title{Interfata grafica realtime: threading, concurenta si afisare}
\author{}
\date{}

\begin{document}

\maketitle
\tableofcontents
\newpage

\section{Introducere}

Interfata grafica realtime a proiectului de Speech Enhancement are rolul de a
demonstra vizual si auditiv functionarea sistemului de reducere a zgomotului.
Aceasta permite captarea semnalului audio de la microfon, procesarea lui prin
metode clasice DSP sau prin modelul neuronal MaskNet, redarea semnalului
imbunatatit catre difuzoare si afisarea in timp real a formelor de unda,
spectrogramelor, nivelurilor audio si scorurilor VAD.

Spre deosebire de un demo offline, interfata realtime trebuie sa execute simultan
mai multe activitati:

\begin{itemize}
    \item citirea continua a semnalului audio de la dispozitivul de intrare;
    \item procesarea fiecarui bloc audio cu latenta mica;
    \item redarea semnalului procesat catre dispozitivul de iesire;
    \item actualizarea grafica a ferestrei;
    \item preluarea comenzilor utilizatorului;
    \item salvarea optionala a inregistrarilor si snapshot-urilor.
\end{itemize}

Pentru a gestiona aceste activitati fara blocaje, proiectul foloseste un model
concurent bazat pe callback audio, thread separat de procesare, cozi pentru
comunicarea intre componente si mecanisme de sincronizare pentru datele partajate.

\section{Modulele implicate}

Interfata realtime este implementata modular, astfel incat partea grafica,
streaming-ul audio si logica de procesare sa fie separate.

\begin{longtable}{p{0.32\textwidth}p{0.60\textwidth}}
\toprule
\textbf{Modul} & \textbf{Rol} \\
\midrule
\endhead
\texttt{src/tools/realtime\_demo.py} &
Punctul de intrare al aplicatiei realtime. Parseaza argumentele CLI, alege
checkpoint-ul, configureaza dispozitivele, creeaza motorul audio si porneste
aplicatia. \\
\addlinespace
\texttt{src/tools/realtime/engine.py} &
Contine clasele \texttt{RealtimeDenoiseApp} si \texttt{StreamingEnhancer}.
Gestioneaza stream-ul audio, thread-ul de procesare, cozile interne si aplicarea
metodelor de denoising. \\
\addlinespace
\texttt{src/tools/realtime/ui.py} &
Contine clasa \texttt{LiveComparisonUI}, responsabila pentru fereastra Tkinter,
controalele interactive, graficele Matplotlib si actualizarea vizuala live. \\
\addlinespace
\texttt{src/tools/realtime/common.py} &
Contine functii si clase comune: alegerea dispozitivelor audio, conversii
milisecunde-esantioane, normalizare, calcul RMS, inregistrare audio si sursa
sintetica \texttt{speech+noise}. \\
\addlinespace
\texttt{src/models/inference.py} &
Contine functiile pentru incarcarea checkpoint-ului MaskNet si aplicarea
modelului asupra unui semnal audio. \\
\bottomrule
\end{longtable}

Aceasta structura permite rularea aceluiasi motor realtime atat cu interfata
grafica, prin \texttt{--visualize}, cat si fara interfata, doar cu loguri in
consola.

\section{Tehnologii folosite pentru concurenta}

\subsection{Callback audio prin sounddevice}

Fluxul audio realtime este gestionat cu biblioteca \texttt{sounddevice}. Aceasta
deschide un stream audio full-duplex, cu un canal de intrare si un canal de
iesire. Stream-ul apeleaza periodic un callback, care primeste un bloc audio nou
si trebuie sa furnizeze rapid un bloc audio pentru iesire.

In proiect, stream-ul este creat cu parametri precum sample rate, dimensiunea
blocului, latenta si dispozitivele active:

\begin{lstlisting}[language=Python]
with sd.Stream(
    samplerate=self.sample_rate,
    blocksize=self.block_size,
    dtype="float32",
    channels=1,
    latency=self.latency,
    device=device,
    callback=self._audio_callback,
):
    ...
\end{lstlisting}

Callback-ul audio este o zona sensibila a aplicatiei. El este apelat foarte des
si nu trebuie sa execute operatii costisitoare, deoarece orice intarziere poate
produce intreruperi audio, dropout-uri sau latenta mare. Din acest motiv,
callback-ul nu ruleaza direct MaskNet si nu face actualizari grafice complexe.
El doar preia intrarea, gestioneaza cozile si trimite la iesire cel mai recent
bloc procesat disponibil.

\subsection{Thread separat pentru procesare}

Procesarea audio este mutata intr-un thread separat, creat cu modulul standard
\texttt{threading}:

\begin{lstlisting}[language=Python]
self.processor_thread = threading.Thread(
    target=self._processor_loop,
    daemon=True
)
\end{lstlisting}

Acest thread ruleaza metoda \texttt{\_processor\_loop}. Rolul lui este sa consume
blocuri audio din coada de intrare, sa aplice metoda de imbunatatire si sa puna
rezultatul in coada de iesire. Astfel, callback-ul audio ramane rapid, iar
operatiile mai lente sunt executate in afara callback-ului.

Separarea este importanta mai ales pentru metoda \texttt{masknet}, unde
procesarea include calcul STFT, inferenta PyTorch si reconstructie prin iSTFT.
Aceste operatii pot dura mai mult decat procesarea clasica DSP si nu trebuie
executate direct in callback-ul stream-ului audio.

\subsection{Model producer-consumer cu queue.Queue}

Comunicarea intre callback-ul audio si thread-ul de procesare se face prin
\texttt{queue.Queue}. Proiectul foloseste doua cozi:

\begin{itemize}
    \item \texttt{input\_queue}, pentru blocurile audio captate;
    \item \texttt{output\_queue}, pentru blocurile audio procesate.
\end{itemize}

Acest mecanism creeaza un model de tip producer-consumer:

\begin{itemize}
    \item callback-ul audio produce blocuri de intrare si le pune in
    \texttt{input\_queue};
    \item thread-ul de procesare consuma blocuri din \texttt{input\_queue};
    \item thread-ul de procesare produce blocuri imbunatatite si le pune in
    \texttt{output\_queue};
    \item callback-ul audio consuma blocuri din \texttt{output\_queue} si le
    trimite catre difuzoare.
\end{itemize}

In forma simplificata, fluxul este:

\begin{lstlisting}[language=Python]
# callback audio
self.input_queue.put_nowait(input_block)
output_block = self.output_queue.get_nowait()

# thread de procesare
block = self.input_queue.get(timeout=0.1)
enhanced = self.enhancer.process_block(block)
self.output_queue.put_nowait(enhanced)
\end{lstlisting}

Coada are dimensiune limitata, configurata prin \texttt{--queue-size}. Daca o
coada se umple, sistemul elimina blocuri vechi sau ignora temporar unele blocuri.
Aceasta alegere favorizeaza continuitatea stream-ului audio in locul acumularii
unei intarzieri din ce in ce mai mari.

\subsection{Oprire controlata cu threading.Event}

Pentru oprirea aplicatiei se foloseste \texttt{threading.Event}. Obiectul
\texttt{stop\_event} este verificat atat in bucla principala, cat si in thread-ul
de procesare:

\begin{lstlisting}[language=Python]
while not self.stop_event.is_set():
    ...
\end{lstlisting}

Cand utilizatorul inchide fereastra, apasa butonul \texttt{Stop Demo} sau
intrerupe executia cu \texttt{Ctrl+C}, evenimentul este setat. Aplicatia opreste
stream-ul audio, asteapta inchiderea thread-ului de procesare si elibereaza
resursele:

\begin{lstlisting}[language=Python]
self.stop_event.set()
self.processor_thread.join(timeout=2.0)
\end{lstlisting}

Aceasta abordare evita oprirea brusca a thread-urilor si permite salvarea
inregistrarilor audio la finalul sesiunii.

\subsection{Sincronizarea datelor cu threading.Lock}

Interfata grafica si callback-ul audio folosesc aceleasi buffere pentru afisarea
semnalului brut si a semnalului imbunatatit. Pentru a evita citirea unor date
modificate partial, clasa \texttt{LiveComparisonUI} foloseste
\texttt{threading.Lock}.

Metoda \texttt{push} actualizeaza bufferele intr-o zona protejata:

\begin{lstlisting}[language=Python]
with self.lock:
    self.raw_buffer = np.roll(self.raw_buffer, -len(raw_block))
    self.raw_buffer[-len(raw_block):] = raw_block
    self.output_buffer = np.roll(self.output_buffer, -len(output_block))
    self.output_buffer[-len(output_block):] = output_block
\end{lstlisting}

Metoda \texttt{refresh} citeste copii ale acestor buffere tot intr-o zona
protejata:

\begin{lstlisting}[language=Python]
with self.lock:
    raw = self.raw_buffer.copy()
    out = self.output_buffer.copy()
\end{lstlisting}

Astfel, UI-ul nu afiseaza date inconsistente, iar callback-ul audio poate continua
sa transmita blocuri noi.

\section{Pipeline-ul concurent de rulare}

Pipeline-ul realtime poate fi descris in pasi succesivi:

\begin{enumerate}
    \item Utilizatorul porneste aplicatia prin CLI:
\begin{lstlisting}[language=bash]
python -m src.tools.project_cli realtime-demo -- --method masknet --visualize
\end{lstlisting}

    \item Modulul \texttt{realtime\_demo.py} citeste argumentele si configureaza
    metoda de denoising, checkpoint-ul, sample rate-ul, dimensiunea blocului,
    dispozitivele audio si optiunile de vizualizare.

    \item Se creeaza obiectul \texttt{StreamingEnhancer}, care pastreaza un
    buffer de context si aplica metoda selectata: \texttt{masknet},
    \texttt{spectral\_subtraction}, \texttt{wiener} sau \texttt{bypass}.

    \item Daca este activata optiunea \texttt{--visualize}, se creeaza fereastra
    \texttt{LiveComparisonUI}.

    \item Se creeaza obiectul \texttt{RealtimeDenoiseApp}, care coordoneaza
    stream-ul audio, thread-ul de procesare, cozile si comunicarea cu UI-ul.

    \item La pornire, motorul face un \texttt{warmup} al procesarii pentru a
    reduce latenta primului bloc.

    \item Este pornit thread-ul de procesare.

    \item Se deschide stream-ul audio \texttt{sounddevice}.

    \item La fiecare callback, aplicatia primeste un bloc audio de intrare.

    \item Blocul de intrare este pus in \texttt{input\_queue}.

    \item Callback-ul incearca sa citeasca un bloc procesat din
    \texttt{output\_queue}. Daca nu exista bloc disponibil, trimite temporar
    liniste la iesire.

    \item Thread-ul de procesare preia blocul din \texttt{input\_queue}, aplica
    denoising si pune rezultatul in \texttt{output\_queue}.

    \item UI-ul primeste blocurile raw si enhanced prin metoda \texttt{push}.

    \item Bucla principala consuma periodic comenzile utilizatorului din UI:
    schimbarea metodei, schimbarea VAD-ului, schimbarea sursei, modificarea
    SNR-ului, snapshot, recording sau oprire.

    \item Metoda \texttt{refresh} actualizeaza graficele si starea vizuala a
    aplicatiei.
\end{enumerate}

\section{Tehnologii folosite pentru afisarea GUI}

\subsection{Tkinter}

Interfata grafica este construita cu \texttt{Tkinter}, toolkit-ul grafic standard
din Python. Fereastra principala este creata in clasa \texttt{LiveComparisonUI}:

\begin{lstlisting}[language=Python]
self.root = tk.Tk()
self.root.title("Realtime Speech Enhancement Monitor")
self.root.geometry("1500x980")
\end{lstlisting}

Tkinter este potrivit pentru acest proiect deoarece este simplu de distribuit,
nu necesita un framework web si permite construirea rapida a unei interfete
desktop cu controale native.

\subsection{ttk pentru widget-uri moderne}

Peste Tkinter este folosit modulul \texttt{ttk}, care ofera widget-uri cu aspect
mai modern si suport pentru stilizare. Interfata foloseste:

\begin{itemize}
    \item \texttt{ttk.Frame} pentru panouri principale;
    \item \texttt{ttk.LabelFrame} pentru sectiuni grupate;
    \item \texttt{ttk.Notebook} pentru taburi;
    \item \texttt{ttk.Combobox} pentru selectii;
    \item \texttt{ttk.Scale} pentru slidere;
    \item \texttt{ttk.Progressbar} pentru niveluri audio;
    \item \texttt{ttk.Button} pentru actiuni precum snapshot, recording si stop.
\end{itemize}

In cod sunt definite stiluri personalizate pentru panouri, butoane, etichete si
bare de progres:

\begin{lstlisting}[language=Python]
style.configure("Card.TFrame", background="#fffdf8")
style.configure("Primary.TButton", font=("Segoe UI", 10, "bold"))
style.configure("Slim.Horizontal.TProgressbar", thickness=10)
\end{lstlisting}

\subsection{Structura ferestrei}

Fereastra este impartita in doua zone:

\begin{itemize}
    \item zona din stanga, pentru monitorizare vizuala;
    \item zona din dreapta, pentru controale interactive.
\end{itemize}

Zona de monitorizare afiseaza:

\begin{itemize}
    \item sumarul sesiunii;
    \item forma de unda a semnalului brut;
    \item forma de unda a semnalului imbunatatit;
    \item spectrograma semnalului brut;
    \item spectrograma semnalului imbunatatit;
    \item nivelurile RMS;
    \item increderea VAD.
\end{itemize}

Zona de control foloseste un \texttt{ttk.Notebook} cu patru taburi:

\begin{itemize}
    \item \textbf{Audio}, pentru sursa audio si dispozitive;
    \item \textbf{Processing}, pentru metoda de denoising, VAD, dry/wet si gain;
    \item \textbf{Files}, pentru incarcarea fisierelor speech/noise;
    \item \textbf{Augment}, pentru scenarii, SNR si distorsiuni.
\end{itemize}

\section{Afisarea graficelor cu Matplotlib}

\subsection{Integrarea Matplotlib in Tkinter}

Graficele live sunt realizate cu \texttt{Matplotlib}. Pentru integrarea figurilor
in fereastra Tkinter se foloseste backend-ul \texttt{FigureCanvasTkAgg}:

\begin{lstlisting}[language=Python]
self.canvas = FigureCanvasTkAgg(fig, master=left)
self.canvas_widget = self.canvas.get_tk_widget()
self.canvas_widget.grid(row=1, column=0, sticky="nsew")
\end{lstlisting}

Figura principala contine patru subplot-uri:

\begin{itemize}
    \item waveform pentru intrarea raw;
    \item waveform pentru iesirea enhanced;
    \item spectrograma pentru intrarea raw;
    \item spectrograma pentru iesirea enhanced.
\end{itemize}

Pentru VAD exista o figura separata, integrata intr-un panou dedicat numit
\texttt{Speech Activity}.

\subsection{Buffere circulare pentru afisare live}

UI-ul pastreaza ultimele esantioane audio in doua buffere NumPy:
\texttt{raw\_buffer} si \texttt{output\_buffer}. Dimensiunea lor este determinata
de parametrul \texttt{display-ms}, care stabileste cate milisecunde de istoric
sunt afisate.

La fiecare bloc nou, bufferele sunt actualizate cu \texttt{np.roll}:

\begin{lstlisting}[language=Python]
self.raw_buffer = np.roll(self.raw_buffer, -len(raw_block))
self.raw_buffer[-len(raw_block):] = raw_block
\end{lstlisting}

Acest mecanism simuleaza o fereastra glisanta peste semnalul audio: datele vechi
sunt impinse spre stanga, iar esantioanele noi apar in partea dreapta a graficului.

\subsection{Actualizarea formelor de unda}

Formele de unda sunt desenate ca obiecte \texttt{Line2D} Matplotlib. La refresh,
nu se recreeaza graficul de la zero, ci se modifica doar datele liniilor:

\begin{lstlisting}[language=Python]
self.raw_line.set_ydata(raw)
self.out_line.set_ydata(out)
\end{lstlisting}

Aceasta abordare este mai eficienta decat redesenarea completa a axelor si ajuta
la mentinerea unei interfete fluide.

\subsection{Calculul spectrogramelor pentru UI}

Spectrogramele afisate in dashboard sunt calculate local in metoda
\texttt{\_compute\_display\_spectrogram}. Pentru fiecare frame audio se aplica o
fereastra Hann, apoi se calculeaza FFT real:

\begin{lstlisting}[language=Python]
fft_frame = np.fft.rfft(frame * window, n=config.N_FFT)
mag_db = 20.0 * np.log10(np.maximum(np.abs(fft_frame), 1e-4))
\end{lstlisting}

Magnitudinea este convertita in decibeli si limitata intre \(-80\) dB si \(0\) dB
pentru afisare stabila. Imaginea spectrogramelor este actualizata prin
\texttt{imshow}:

\begin{lstlisting}[language=Python]
self.raw_spec_image.set_data(raw_spec)
self.out_spec_image.set_data(out_spec)
\end{lstlisting}

Spectrograma raw foloseste colormap-ul \texttt{magma}, iar spectrograma enhanced
foloseste \texttt{viridis}, pentru diferentiere vizuala rapida.

\subsection{Afisarea VAD live}

Interfata afiseaza si o curba de incredere VAD. La fiecare refresh, pe bufferul
raw este rulat detectorul \texttt{EnergyZCRVAD} in modul soft:

\begin{lstlisting}[language=Python]
vad_soft = EnergyZCRVAD().predict_soft(raw)
self.vad_line.set_data(vad_time, vad_soft)
\end{lstlisting}

Curba rezultata are valori in intervalul \([0,1]\), unde valorile apropiate de
1 indica probabilitate mai mare de activitate vocala.

\subsection{Limitarea ratei de redesenare}

Redesenarea grafica este costisitoare, mai ales cand sunt actualizate simultan
waveform-uri si spectrograme. Pentru a evita consumul excesiv de CPU, canvas-ul
este redesenat doar daca a trecut un interval minim de timp:

\begin{lstlisting}[language=Python]
if now - self.last_refresh_time >= 0.08:
    self.canvas.draw_idle()
    self.last_refresh_time = now
\end{lstlisting}

Aceasta inseamna o frecventa vizuala de aproximativ 12 cadre pe secunda, suficienta
pentru monitorizare audio, dar mult mai eficienta decat redesenarea la fiecare
callback audio.

\section{Gestionarea comenzilor din interfata}

\subsection{Variabile pending}

Comenzile utilizatorului nu modifica direct starea motorului audio. In schimb,
interfata seteaza variabile intermediare de tip \texttt{pending}. De exemplu,
cand utilizatorul schimba metoda de denoising, UI-ul actualizeaza
\texttt{pending\_method}:

\begin{lstlisting}[language=Python]
def _on_method_var_changed(self, *_args) -> None:
    self.pending_method = self.method_var.get()
\end{lstlisting}

Motorul realtime verifica periodic daca exista o cerere noua:

\begin{lstlisting}[language=Python]
requested_method = self.visualizer.consume_requested_method()
if requested_method is not None:
    self.enhancer.set_method(requested_method)
\end{lstlisting}

Acest model este folosit pentru metoda de procesare, modul VAD, sursa audio,
dispozitive, SNR, dry/wet, output gain, scenarii, distorsiuni, snapshot si
recording.

\subsection{Avantajul modelului pending-consume}

Modelul \texttt{pending-consume} are doua avantaje principale. In primul rand,
evita modificarea directa a motorului audio din callback-urile GUI. In al doilea
rand, toate schimbarile sunt aplicate intr-un punct controlat al buclei realtime,
ceea ce reduce riscul de conditii de cursa si stari partial actualizate.

\section{Robustete si tratarea problemelor realtime}

Aplicatia include mai multe mecanisme de robustete:

\begin{itemize}
    \item daca \texttt{output\_queue} este goala, se trimite temporar liniste la
    iesire;
    \item daca \texttt{input\_queue} este plina, se elimina blocuri vechi pentru
    a evita cresterea latentei;
    \item daca microfonul pare silentios pentru mai multe callback-uri, se emite
    un avertisment;
    \item daca schimbarea dispozitivelor audio esueaza, aplicatia revine la
    ultima pereche de dispozitive functionala;
    \item la inchidere, thread-ul de procesare este oprit controlat;
    \item daca recording-ul este activ, fisierele raw si enhanced sunt salvate
    la final.
\end{itemize}

\section{Concluzie}

Interfata grafica realtime foloseste o arhitectura concurenta simpla, dar
potrivita pentru procesarea audio interactiva. Callback-ul \texttt{sounddevice}
este mentinut scurt, procesarea este mutata intr-un thread separat, comunicarea
se face prin cozi, oprirea este controlata prin \texttt{threading.Event}, iar
bufferele partajate sunt protejate prin \texttt{threading.Lock}.

Pentru afisare, proiectul foloseste \texttt{Tkinter} si \texttt{ttk} pentru
fereastra si controale, iar \texttt{Matplotlib} pentru forme de unda, spectrograme
si curba VAD. Prin combinarea acestor tehnologii, aplicatia poate rula in timp
real, poate raspunde la comenzile utilizatorului si poate afisa clar diferenta
dintre semnalul zgomotos si semnalul imbunatatit.

\end{document}
